% Generated draft: requires mathematical and source review.
\documentclass[11pt]{article}
\usepackage[margin=25mm]{geometry}
\usepackage{fontspec}
\setmainfont{Latin Modern Roman}
\usepackage{amsmath,amssymb,mathtools,unicode-math}
\setmathfont{Latin Modern Math}
\usepackage{enumitem,needspace,xurl,hyperref}
\setlength{\parindent}{0pt}
\setlength{\parskip}{5pt}
\setlength{\emergencystretch}{3em}
\newcommand{\N}{\mathbb{N}}
\newcommand{\Z}{\mathbb{Z}}
\newcommand{\Q}{\mathbb{Q}}
\newcommand{\R}{\mathbb{R}}
\newcommand{\C}{\mathbb{C}}
\newcommand{\F}{\mathbb{F}}
\newcommand{\A}{\mathbb{A}}
\newcommand{\PP}{\mathbb{P}}
\newcommand{\E}{\mathbb{E}}
\newcommand{\eps}{\varepsilon}
\newcommand{\dd}{\,\mathrm{d}}
\begin{document}
\section*{UnsolvedMath ID: 800179}
\textbf{Source problem code:} catalogue entry
\subsection{Definitions and mathematical statement}
For each integer \(k\ge 3\), write formally
\[
E_{k,2}(x,h):=D_{k,2}(x,h)-M_{k,2}(x,h),
\]
where the supplied source describes \(M_{k,2}(x,h)\) as the main term in an asymptotic formula for \(D_{k,2}(x,h)\).

\textbf{Requested statement (N. Ng, Lethbridge).} Prove an omega theorem for \(E_{k,2}(x,h)\) for \(k\ge 3\).

\textbf{Unresolved qualifications.} The supplied record does not define \(D_{k,2}(x,h)\), specify the main term \(M_{k,2}(x,h)\), give the domains or ranges of \(x\) and \(h\), or state whether \(h\) is fixed or varies with \(x\). It also does not specify an omega scale or whether an unsigned or signed omega result is intended. Thus the request cannot be converted into a uniquely specified mathematical proposition from this record alone.

For clarity, a possible unsigned omega assertion, once the missing data are supplied, has the following meaning. For an unbounded set \(X\subseteq\R\), functions \(E:X\to\C\) and \(B:X\to(0,\infty)\), the assertion
\[
E(x)=\Omega(B(x))\qquad(x\to\infty,\ x\in X)
\]
means that there exist \(c>0\) and a sequence \((x_j)_{j\ge1}\) in \(X\), tending to infinity, such that
\[
|E(x_j)|\ge cB(x_j)\qquad\text{for every }j\ge1.
\]
This definition is an explanatory template, not a selection of the intended domain, scale, sign convention, or parameter quantifiers for the source problem. Any dependence of constants and sequences on \(k\) and \(h\), and any required uniformity, remains unspecified.
\subsection{Short English statement}
N. Ng asks for an omega theorem for the error after subtracting the main term from \(D_{k,2}(x,h)\), for integers \(k\ge3\). The supplied record omits the definitions, parameter ranges, main term, and intended omega scale and sign, so a complete precise formulation requires those details.
\subsection{Sources}
\begin{itemize}
\item Local UnsolvedMath registry, permanent ID \(800179\), exact source code \texttt{catalogue entry}, supplied complete source record; upstream revision \texttt{d371ee917ac9647fbb06cd6fc529deec6f3dda82}. The registry labels this problem \texttt{open}; this is a supplied status label, not an independently verified mathematical assessment.
\item Hamieh, Kadiri, Ng, \emph{Moments of L-functions Problem List} (2024), Problem 30, in the supplied category ``Moments of the Riemann zeta function'': \url{https://arxiv.org/abs/2405.17800}. The supplied provenance records access on 2026-10-16 and extraction by \texttt{arxiv\_latex}; these provenance details have not been independently verified here.
\item The problem text attributes the request to N. Ng, Lethbridge. The registry records Hamieh, Kadiri, Ng as proposers and 2024 as the proposed year.
\item The supplied statement credits Topacogullari, references [T1] and [T3], and Drappeau, reference [Dr], with asymptotic formulae for \(D_{k,2}(x,h)\). This is a report of the source's background claim. Bibliographic entries, theorem statements, hypotheses, and links for these three reference labels are not supplied.
\end{itemize}
\end{document}
